\section*{Hoofdstuk \ref{ch:b1:functions} --- Standaardfuncties}

\begin{solution}{exo:b1:functions:1}
$\arcsin\frac{\sqrt 3}{2} = \frac\pi3$;
$\;\arccos\bigl(-\frac12\bigr) = \frac{2\pi}{3}$;
$\;\arctan(-1) = -\frac\pi4$.

$\arcsin\bigl(\sin\frac{5\pi}{6}\bigr)$: er is $\sin\frac{5\pi}{6} =
\frac12$, en de hoek uit $\intcc{-\frac\pi2}{\frac\pi2}$ met sinus
$\frac12$ is $\frac\pi6$ (niet $\frac{5\pi}{6}$).

$\arccos\bigl(\cos\frac{7\pi}{4}\bigr)$: er is $\cos\frac{7\pi}{4} =
\frac{\sqrt 2}{2}$, en de hoek uit $\intcc{0}{\pi}$ met die cosinus is
$\frac\pi4$.
\end{solution}

\begin{solution}{exo:b1:functions:2}
$f(x) = \arcsin(2x - 1)$ vergt $-1 \leq 2x - 1 \leq 1$, oftewel $x \in
\intcc{0}{1}$. Op $\intoo{0}{1}$ geven de kettingregel en
\cref{prop:b1:functions:arcderiv}
\[
f'(x) = \frac{2}{\sqrt{1 - (2x-1)^2}} = \frac{2}{\sqrt{4x - 4x^2}}
= \frac{1}{\sqrt{x(1 - x)}} .
\]
$g(x) = \arctan\sqrt x$ is gedefinieerd voor $x \geq 0$, en voor $x > 0$
is
\[
g'(x) = \frac{1}{1 + x} \cdot \frac{1}{2\sqrt x}
= \frac{1}{2\sqrt x\,(1 + x)} .
\]
\end{solution}

\begin{solution}{exo:b1:functions:3}
$\cosh x + \sinh x = \frac{\eu^x + \eu^{-x}}{2} + \frac{\eu^x -
\eu^{-x}}{2} = \eu^x$. Bijgevolg is voor $n \in \Z$
\[
(\cosh x + \sinh x)^n = (\eu^x)^n = \eu^{nx}
= \cosh nx + \sinh nx .
\]
\end{solution}

\begin{solution}{exo:b1:functions:4}
Met $y = \arcsin x$: $\cos 2y = 1 - 2\sin^2 y = 1 - 2x^2$, dus
$\cos(2\arcsin x) = 1 - 2x^2$.

Met $y = \arctan x$: $\sin 2y = 2 \sin y \cos y = 2 \tan y \cos^2 y =
\frac{2\tan y}{1 + \tan^2 y}$, dus $\sin(2\arctan x) =
\dfrac{2x}{1 + x^2}$.
\end{solution}

\begin{solution}{exo:b1:functions:5}
Van traagst naar snelst:
\[
(\ln x)^{1000} \;\ll\; x^{100} \;\ll\; x^{\ln x} \;\ll\;
\eu^{\sqrt x} \;\ll\; \eu^{x/100} .
\]
Verantwoording: $(\ln x)^{1000}/x^{100} \to 0$ volgens
\cref{prop:b1:functions:powerrules} (logaritmen verliezen van machten).
Verder is $x^{100} = \eu^{100\ln x}$ en $x^{\ln x} = \eu^{(\ln x)^2}$:
omdat $(\ln x)^2 - 100\ln x \to +\infty$, wint de tweede. Vervolgens is
$x^{\ln x} = \eu^{(\ln x)^2} \ll \eu^{\sqrt x}$ omdat $(\ln x)^2/\sqrt x
\to 0$ (logaritmen verliezen van de macht $x^{1/4}$, gekwadrateerd). Ten
slotte gaat $\sqrt x - x/100 \to -\infty$, dus $\eu^{\sqrt x} \ll
\eu^{x/100}$.
\end{solution}

\begin{solution}{exo:b1:functions:6}
Domein: $x \neq \pm 1$, drie intervallen. Op elk ervan is
\[
f'(x) = \frac{\bigl(\frac{2x}{1-x^2}\bigr)'}{1 +
\frac{4x^2}{(1-x^2)^2}}
= \frac{\frac{2(1-x^2) + 4x^2}{(1-x^2)^2}}
{\frac{(1-x^2)^2 + 4x^2}{(1-x^2)^2}}
= \frac{2(1 + x^2)}{(1 + x^2)^2} = \frac{2}{1 + x^2}
= (2\arctan x)' .
\]
Dus is $f(x) - 2\arctan x$ constant op elk interval. Waarden: in $x = 0$
is $f(0) = 0$, zodat de constante $0$ is op $\intoo{-1}{1}$. Voor $x \to
+\infty$ gaat $\frac{2x}{1-x^2} \to 0^-$, dus $f \to 0$, terwijl
$2\arctan x \to \pi$: de constante is $-\pi$ op $\intoo{1}{+\infty}$.
Wegens de oneven pariteit is ze $+\pi$ op $\intoo{-\infty}{-1}$.
Samengevat: $f = 2\arctan x$ op $\intoo{-1}{1}$, $f = 2\arctan x - \pi$
voor $x > 1$, en $f = 2\arctan x + \pi$ voor $x < -1$. (Dit is de
dubbelehoekformule voor de tangens, gelezen door $\arctan$ heen.)
\end{solution}

\begin{solution}{exo:b1:functions:7}
$\cosh x = 2$: met $u = \eu^x > 0$ wordt dat $u + u^{-1} = 4$, dus $u^2
- 4u + 1 = 0$ en $u = 2 \pm \sqrt 3$: $x = \ln(2 + \sqrt 3)$ of $x =
\ln(2 - \sqrt 3) = -\ln(2 + \sqrt 3)$ (de twee oplossingen zijn
elkaars tegengestelde, want $\cosh$ is even; beide zijn geldig).
Equivalent: $x = \pm \operatorname{arcosh} 2$.

$5\cosh x - 4\sinh x = 3$: na substitutie van de exponentiële definities
wordt dat $\frac{5(u + u^{-1}) - 4(u - u^{-1})}{2} = 3$, oftewel $u +
9u^{-1} = 6$, oftewel $u^2 - 6u + 9 = (u - 3)^2 = 0$: $u = 3$, met de
ene oplossing $x = \ln 3$.
\end{solution}

\begin{solution}{exo:b1:functions:8}
Zij $\alpha = \arctan\frac12 + \arctan\frac13$. De somformule geeft
\[
\tan\alpha = \frac{\frac12 + \frac13}{1 - \frac12\cdot\frac13}
= \frac{5/6}{5/6} = 1 .
\]
Beide boogtangenten liggen in $\intoo{0}{\frac\pi4}$ (hun argumenten
liggen in $\intoo{0}{1}$), dus $\alpha \in \intoo{0}{\frac\pi2}$; de
enige hoek daar met tangens $1$ is $\frac\pi4$.

Machin: zij $\beta = \arctan\frac15$. Tweemaal de dubbelehoekformule:
\[
\tan 2\beta = \frac{2/5}{1 - 1/25} = \frac{5}{12},
\qquad
\tan 4\beta = \frac{2 \cdot 5/12}{1 - 25/144} = \frac{120}{119}.
\]
Vervolgens is
\[
\tan\Bigl(4\beta - \frac\pi4\Bigr)
= \frac{\frac{120}{119} - 1}{1 + \frac{120}{119}}
= \frac{1/119}{239/119} = \frac{1}{239}.
\]
Lokalisering: $\beta < \arctan 1 = \frac\pi4$, en zelfs $4\beta \in
\intoo{0}{\frac\pi2}$ met $\tan 4\beta = \frac{120}{119}$ dicht bij $1$
(want $\beta < \frac\pi8$, aangezien $\tan\frac\pi8 = \sqrt 2 - 1 >
\frac15$); dus ligt $4\beta - \frac\pi4$ in
$\intoo{-\frac\pi4}{\frac\pi4}$, waar $\arctan$ de tangens inverteert:
$4\beta - \frac\pi4 = \arctan\frac{1}{239}$, en dat is de formule van
Machin.
\end{solution}

\begin{solution}{exo:b1:functions:9}
Zij $f(x) = x - \sin x$: dan is $f(0) = 0$ en $f'(x) = 1 - \cos x \geq
0$, dus $f \geq 0$ op $\R_+$: $\sin x \leq x$.

Zij $g(x) = \sin x - x + \frac{x^3}{6}$: dan is $g(0) = 0$, $g'(x) =
\cos x - 1 + \frac{x^2}{2}$, $g'(0) = 0$ en $g''(x) = -\sin x + x = f(x)
\geq 0$ op $\R_+$. Dus is $g'$ stijgend met $g'(0) = 0$, waardoor $g'
\geq 0$, waardoor $g$ stijgend is met $g(0) = 0$: $g \geq 0$, oftewel $x
- \frac{x^3}{6} \leq \sin x$ op $\R_+$.

Bijgevolg is $0 \leq \frac{x - \sin x}{x^3} \leq \frac 16$ voor $x > 0$:
de verhouding zit opgesloten op de schaal van $\frac16$, en
\cref{ch:b1:taylor} toont aan dat haar limiet precies $\frac 16$ is.
\end{solution}

\begin{solution}{exo:b1:functions:10}
Schrijf $y = \arcsin x$, zodat $2x\sqrt{1 - x^2} = 2\sin y\cos y = \sin
2y$ en $h(x) = \arcsin(\sin 2y) - 2y$.

Voor $\abs x \leq \frac{\sqrt 2}{2}$: dan is $2y \in
\intcc{-\frac\pi2}{\frac\pi2}$, dus $\arcsin(\sin 2y) = 2y$ en $h = 0$.

Voor $x \in \intoc{\frac{\sqrt 2}{2}}{1}$: dan is $2y \in
\intoc{\frac\pi2}{\pi}$, en de hoek uit $\intcc{-\frac\pi2}{\frac\pi2}$
met sinus $\sin 2y$ is $\pi - 2y$: dus $h(x) = \pi - 4\arcsin x$.
(Controle via de afgeleide: daar is $h'(x) = \frac{2(1 - 2x^2)}{\abs{1 -
2x^2}\sqrt{1 - x^2}} - \frac{2}{\sqrt{1-x^2}} = \frac{-4}{\sqrt{1 -
x^2}}$, de afgeleide van $-4\arcsin x$; en in $x = 1$ is $h(1) = \arcsin
0 - 2\cdot\frac\pi2 = -\pi = \pi - 4\cdot\frac\pi2$.)

Voor $x \in \intco{-1}{-\frac{\sqrt 2}{2}}$ geeft de oneven pariteit van
$h$: $h(x) = -\pi - 4\arcsin x$.
\end{solution}

\begin{solution}{exo:b1:functions:11}
$g = \arctan \circ \sinh$ is een samenstelling van oneven, strikt
stijgende functies, dus zelf oneven en strikt stijgend; voor $x \to
+\infty$ gaat $\sinh x \to +\infty$, dus $g(x) \to \frac\pi2$, en $g$ is
een bijectie op $\intoo{-\frac\pi2}{\frac\pi2}$ (continuïteit plus de
tussenwaardestelling). De kettingregel met
\cref{prop:b1:functions:hyprules}~(1) geeft
\[
g'(x) = \frac{\cosh x}{1 + \sinh^2 x} = \frac{\cosh x}{\cosh^2 x}
= \frac{1}{\cosh x} .
\]
Per constructie is $\tan g(x) = \sinh x$. Omdat $g(x) \in
\intoo{-\frac\pi2}{\frac\pi2}$ een positieve cosinus heeft, volgt
\begin{align*}
\cos g(x) &= \frac{1}{\sqrt{1 + \tan^2 g(x)}}
= \frac{1}{\sqrt{1 + \sinh^2 x}} = \frac{1}{\cosh x},\\
\sin g(x) &= \tan g(x) \cos g(x) = \frac{\sinh x}{\cosh x} = \tanh x .
\end{align*}
\end{solution}

\begin{solution}{exo:b1:functions:12}
Zet $u = \arctan 2$ en $v = \arctan 3$; beide liggen in
$\intoo{\frac\pi4}{\frac\pi2}$ (hun argumenten zijn groter dan $1$),
dus $u + v \in \intoo{\frac\pi2}{\pi}$. De somformule geeft
\[
\tan(u + v) = \frac{2 + 3}{1 - 6} = -1 ,
\]
en de unieke hoek uit $\intoo{\frac\pi2}{\pi}$ met tangens $-1$ is
$\frac{3\pi}4$: dus $\arctan 2 + \arctan 3 = \frac{3\pi}4$. (Blindelings
$\arctan$ op de tangens toepassen zou $-\frac\pi4$ geven, een $\pi$
ernaast --- de lokalisering redt de berekening, en de negatieve noemer
$1 - 6$ is juist het signaal dat de som
$\intoo{-\frac\pi2}{\frac\pi2}$ verlaten heeft.) Optellen van $\arctan 1
= \frac\pi4$ geeft
\[
\arctan 1 + \arctan 2 + \arctan 3 = \frac\pi4 + \frac{3\pi}4 = \pi .
\]
\end{solution}

\begin{solution}{pb:b1:functions:1}
\textbf{1.} $f$ is een quotiënt van afleidbare functies met een noemer
die niet nul wordt op $\intoo0{+\infty}$, en
\[
f'(t) = \frac{\frac1t \cdot t - \ln t}{t^2} = \frac{1 - \ln
t}{t^2} ,
\]
positief voor $t < \eu$, nul in $t = \eu$, negatief voor $t > \eu$: $f$
stijgt strikt op $\intoc0\eu$ en daalt strikt op $\intco\eu{+\infty}$,
met maximum $f(\eu) = \frac1\eu$.

\textbf{2.} Voor $t \to 0^+$ gaat $\ln t \to -\infty$ en $\frac1t \to
+\infty$, dus $f(t) \to -\infty$. Voor $t \to +\infty$ gaat $f(t) \to 0$
wegens de \omterm{prop:b1:functions:powerrules}{groeivergelijking} van \cref{prop:b1:functions:powerrules}
($\beta = \alpha = 1$). Teken: dat van $\ln t$, dus $f < 0$ op
$\intoo01$, $f(1) = 0$ en $f > 0$ op $\intoo1{+\infty}$. De grafiek
klimt vanuit $-\infty$, snijdt de as in $1$, piekt in $(\eu,
\frac1\eu)$ en zakt daarna naar $0^+$.

\textbf{3.} $f(4) = \frac{\ln 4}4 = \frac{2\ln 2}4 = \frac{\ln 2}2 =
f(2)$.

\textbf{4.} Met de verlooptabel en de tussenwaardestelling (hier
gebruikt zoals bekend uit het bovenbouwvolume; geformaliseerd in
\cref{ch:b1:continuity}). Voor $c < 0$ bestaan er alleen oplossingen waar
$f < 0$, dat wil zeggen in $\intoo01$, waar $f$ een strikt stijgende
bijectie op $\intoo{-\infty}0$ is: precies één. Voor $c = 0$: alleen
$t = 1$. Voor $0 < c < \frac1\eu$: op $\intoo1\eu$ stijgt $f$ van $0$
naar $\frac1\eu$, wat één oplossing geeft; op $\intoo\eu{+\infty}$ daalt
$f$ van $\frac1\eu$ naar $0$, wat er nog één geeft; samen twee. Voor $c
= \frac1\eu$: alleen het maximumpunt $t = \eu$. Voor $c > \frac1\eu$:
geen.

\textbf{5.} Voor $x \neq \eu$ geeft de strikte maximaliteit $f(x) <
f(\eu)$, oftewel $\frac{\ln x}x < \frac1\eu$, oftewel $\eu\ln x < x$,
oftewel $\ln(x^\eu) < \ln(\eu^x)$: $x^\eu < \eu^x$. Met $x = \pi$ volgt
$\eu^\pi > \pi^\eu$ (numeriek $23.14 > 22.46$).

\textbf{6.} Voor $x, y > 0$ zijn beide leden positief, dus
\[
x^y = y^x \iff y\ln x = x\ln y \iff \frac{\ln x}x = \frac{\ln y}y
\iff f(x) = f(y) ,
\]
na deling door $xy > 0$.

\textbf{7.} Zij $f(x) = f(y) = c$ met $x \neq y$. Is $c \leq 0$ of $c =
\frac1\eu$, dan zegt vraag 4 dat de vergelijking $f = c$ één enkele
oplossing heeft: onmogelijk. Dus $0 < c < \frac1\eu$, en opnieuw volgens
vraag 4 zijn de twee oplossingen één punt van $\intoo1\eu$ en één van
$\intoo\eu{+\infty}$: beide coördinaten zijn groter dan $1$ en ze
liggen aan weerszijden van $\eu$.

\textbf{8.} Substitutie van $y = tx$ in $y\ln x = x\ln y$ geeft
\[
tx\ln x = x(\ln t + \ln x)
\iff (t - 1)\ln x = \ln t
\iff \ln x = \frac{\ln t}{t - 1} ,
\]
dus $x = t^{1/(t-1)}$ en $y = tx = t^{1 + \frac1{t-1}} = t^{t/(t-1)}$.
Omgekeerd is voor die waarden $\ln y = \frac{t\ln t}{t-1} = t\ln x$ en $y
= tx$, zodat $\frac{\ln y}y = \frac{t\ln x}{tx} = \frac{\ln x}x$: een
oplossing. Elke niet-triviale oplossing heeft een zekere verhouding $t =
y/x \in \intoo0{+\infty} \setminus \{1\}$, zodat de parametrisering
volledig is.

\textbf{9.} $t = 2$: $x = 2^{1/1} = 2$ en $y = 2^{2/1} = 4$. En
\[
x(1/t) = (1/t)^{\frac1{\frac1t - 1}}
= (t^{-1})^{\frac{t}{1 - t}}
= t^{\frac{t}{t - 1}} = y(t) ,
\]
waarna $y(1/t) = \frac1t\,x(1/t) = \frac{y(t)}t = x(t)$: de
parameterwissel $t \mapsto 1/t$ verwisselt de coördinaten, zoals de
symmetrie van de vergelijking eist.

\textbf{10.} Voor $t \to 1$ gaat $\frac{\ln t}{t-1} \to 1$ (het is het
differentiequotiënt van $\ln$ in $1$), dus $x \to \eu^1 = \eu$ en $y =
tx \to \eu$. Voor $t \to +\infty$ gaat $\ln x = \frac{\ln t}{t-1} \to 0$,
dus $x \to 1$, terwijl $\ln y = \frac{t\ln t}{t-1} \to +\infty$, dus $y
\to +\infty$. Voor $t \to 0^+$ gaat $\ln x = \frac{\ln t}{t-1} \to
\frac{-\infty}{-1} = +\infty$, dus $x \to +\infty$, terwijl $\ln y =
\frac{t\ln t}{t-1} \to \frac{0}{-1} = 0$, dus $y \to 1$ (met $t\ln t \to
0$, \cref{ex:b1:functions:xx}). De tak loopt dus van de asymptoot $y = 1$
(uiterst rechts) omhoog door $(\eu, \eu)$ op de diagonaal en weg langs de
asymptoot $x = 1$ (uiterst boven) --- symmetrisch om de diagonaal wegens
vraag 9.

\textbf{11.} Er is $\ln x(t) = g(t) = \frac{\ln t}{t-1}$ en
\[
g'(t) = \frac{\frac{t-1}t - \ln t}{(t-1)^2}
= \frac{h(t)}{(t-1)^2},
\qquad h(t) = 1 - \frac1t - \ln t .
\]
Er geldt $h(1) = 0$ en $h'(t) = \frac1{t^2} - \frac1t = \frac{1 -
t}{t^2} < 0$ voor $t > 1$: dus $h < 0$ op $\intoo1{+\infty}$, waardoor
$g' < 0$ en $x = \eu^g$ daar strikt dalend is (van $\eu$ naar $1$,
wegens vraag 10).

\textbf{12.} Volgens vraag 11 is $t \mapsto x(t)$ een strikt dalende
bijectie van $\intoo1{+\infty}$ op $\intoo1\eu$; noteer $t(x)$ voor haar
inverse (eveneens strikt dalend) en zet $\varphi(x) = y(t(x))$. Merk op
dat $h < 0$ ook op $\intoo01$ (daar is $h' > 0$ en $h(1) = 0$), zodat
$x(\cdot)$ ook op $\intoo01$ dalend is; bijgevolg is $y(t) = x(1/t)$
(vraag 9) \emph{stijgend} in $t$ op $\intoo1{+\infty}$, van $\eu$ naar
$+\infty$. Samenstellen geeft: $\varphi = y \circ t(\cdot)$ is strikt
dalend van $\intoo1\eu$ op $\intoo\eu{+\infty}$, en $x^{\varphi(x)} =
\varphi(x)^x$ volgens vraag 8. Ten slotte is voor een gemeenschappelijke
waarde $c = f(x) \in \intoo0{\frac1\eu}$ het paar $\{x, \varphi(x)\}$ dé
oplossingsverzameling met twee elementen van $f = c$ (vraag 4); breid je
$\varphi$ uit tot $\intoo\eu{+\infty}$ als de inverse \omterm{def:b1:logic:map}{afbeelding}, dan
keert $\varphi(\varphi(x))$ terug naar het andere (dus oorspronkelijke)
element: $\varphi \circ \varphi = \mathrm{id}$.

\textbf{13.} Is $(x, y)$ een niet-triviale geheeltallige oplossing met
$x < y$, dan legt vraag 7 $x$ in $\intoo1\eu$: het enige gehele getal
daar is $x = 2$. Dan is $f(y) = f(2) = \frac{\ln2}2$ met $y > \eu$;
volgens vraag 4 heeft de vergelijking $f = \frac{\ln2}2$ precies één
oplossing voorbij $\eu$, en vraag 3 wijst haar aan: $y = 4$. Dus zijn
$(2, 4)$ en zijn verwisseling de enige.

\textbf{14.} $t = 1 + \frac1n = \frac{n+1}n$ geeft $\frac1{t - 1} = n$
en $\frac t{t-1} = n + 1$, dus
\[
x_n = \Bigl(\frac{n+1}n\Bigr)^{n},
\qquad
y_n = \Bigl(\frac{n+1}n\Bigr)^{n+1},
\]
klaarblijkelijk rationaal en verschillend ($y_n = t\,x_n \neq x_n$), en
$n = 1$ geeft $x_1 = 2$ en $y_1 = 4$.

\textbf{15.} $t = y/x$ is rationaal en $> 1$; schrijf $t - 1 = \frac rs$
onvereenvoudigbaar, zodat $t = \frac{s + r}s$ met $\gcd(s + r, s) =
\gcd(r, s) = 1$. Vraag 8 geeft $x = t^{s/r}$, dus
\[
x^r = t^s = \frac{(s+r)^s}{s^s} ,
\]
een onvereenvoudigbare breuk (geen enkel priemgetal deelt zowel $s + r$
als $s$). Schrijven we $x = \frac pq$ onvereenvoudigbaar, dan is $x^r =
\frac{p^r}{q^r}$ eveneens onvereenvoudigbaar, en wegens de uniciteit van
die schrijfwijze is $p^r = (s+r)^s$ en $q^r = s^s$. Vergelijk de exponent
van een willekeurig priemgetal $\ell$ in $q^r = s^s$: $r \cdot v_\ell(q)
= s \cdot v_\ell(s)$, dus $r$ deelt $s\,v_\ell(s)$; omdat $\gcd(r, s) =
1$, deelt $r$ het getal $v_\ell(s)$ voor elk priemgetal $\ell$
(uniciteit van de priemontbinding, aangenomen; bewezen in
\cref{ch:b1:arith}), zodat $s = b^r$ voor een geheel getal $b$.
Hetzelfde argument op $p^r = (s+r)^s$ geeft $s + r = a^r$.

\textbf{16.} Stel $a^r - b^r = r$ met gehele getallen $a > b \geq 1$ en
$r \geq 2$. Dan is $a \geq b + 1$, zodat het binomium geeft
\[
a^r - b^r \geq (b+1)^r - b^r
= \sum_{k=0}^{r-1}\binom rk b^k
\geq \binom r0 + \binom r{r-1} b^{r-1}
\geq 1 + r > r ,
\]
een tegenspraak. Volgens vraag 15 dwingt $a^r - b^r = (s + r) - s = r$
dus af dat $r = 1$: $t = 1 + \frac1s$, en de oplossing is het paar
$(x_s, y_s)$ van vraag 14. Samen met de verwisselde paren is de
classificatie volledig.

\textbf{17.} $f\bigl(\tfrac94\bigr) = \frac{\ln 2.25}{2.25} =
\frac{0.81093}{2.25} = 0.3604$ en $f\bigl(\tfrac{27}8\bigr) = \frac{\ln
3.375}{3.375} = \frac{1.21640}{3.375} = 0.3604$ (vier decimalen):
gelijk, zoals de constructie belooft --- dus
$\bigl(\tfrac94\bigr)^{27/8} = \bigl(\tfrac{27}8\bigr)^{9/4}$.

\textbf{18.} De parameters $t_n = 1 + \frac1n$ dalen strikt naar $1$.
Omdat $x(\cdot)$ strikt dalend is op $\intoo1{+\infty}$ (vraag 11), is
$x_n = x(t_n)$ strikt \emph{stijgend}; omdat $y(\cdot)$ daar strikt
stijgend is (vraag 12), is $y_n = y(t_n)$ strikt \emph{dalend}. Volgens
vraag 7 is $x_n \in \intoo1\eu$ en $y_n \in \intoo\eu{+\infty}$: dus
$x_n < \eu < y_n$ voor elke $n$. Ten slotte geven $t_n \to 1$ en vraag
10 dat $x_n \to \eu$ en $y_n \to \eu$. De klassieke monotone convergentie
van $\bigl(1 + \frac1n\bigr)^n$ valt zo uit de meetkunde van de tak, en
er kwam geen enkele nieuwe ongelijkheid aan te pas.

\textbf{19.} Is $\eu \leq a < b$, dan geeft $f$ strikt dalend op
$\intco\eu{+\infty}$ dat $f(a) > f(b)$, oftewel $b\ln a > a\ln b$,
oftewel $a^b > b^a$. Is $1 < a < b \leq \eu$, dan geeft $f$ strikt
stijgend dat $f(a) < f(b)$, dus $a^b < b^a$. Gemengd geval: $2 < \eu <
3$ met $2^3 = 8 < 9 = 3^2$, maar $2 < \eu < 5$ met $2^5 = 32 > 25 =
5^2$: beide uitkomsten komen voor, beslist door aan welke kant van de tak
het punt $(a, b)$ valt.

\textbf{20.} De tak snijdt de diagonaal in $\eu$ en $\varphi$ zet zich
daar continu voort met $\varphi(\eu) = \eu$. Differentiëren van
$\varphi(\varphi(x)) = x$ in $x = \eu$ met de kettingregel geeft
$\varphi'(\varphi(\eu))\,\varphi'(\eu) = \varphi'(\eu)^2 = 1$, dus
$\varphi'(\eu) = \pm1$; omdat $\varphi$ dalend is, is $\varphi'(\eu) =
-1$. De tak snijdt de diagonaal met richtingscoëfficiënt $-1$: loodrecht.

\textbf{21.} $y = 3x$ is het geval $t = 3$: $x = 3^{1/2} = \sqrt3$ en $y
= 3^{3/2} = 3\sqrt3$. Controle: $f(\sqrt3) = \frac{0.54931}{1.73205} =
0.3171$ en $f(3\sqrt3) = \frac{\ln 5.19615}{5.19615} =
\frac{1.64792}{5.19615} = 0.3171$: gelijk tot op vier decimalen, dus
$(\sqrt3)^{3\sqrt3} = (3\sqrt3)^{\sqrt3}$.

\textbf{22.} Er is $n^{1/n} = \eu^{f(n)}$, en $\exp$ is stijgend, zodat
de rangschikking die van $f(n)$ is: $f(3) = \frac{\ln3}3 \approx 0.3662$
overtreft $f(2) = f(4) = \frac{\ln2}2 \approx 0.3466$ (vraag 3) en $f(5)
\approx 0.3219$. De grootste is $\sqrt[3]3$, en het gelijke paar is
$\sqrt2 = \sqrt[4]4$ --- het geheeltallige paar $(2, 4)$ in weer een
andere vermomming.

\textbf{23.} De oplossingsverzameling is de vereniging van de diagonaal
$\{(x, x) : x > 0\}$ en één enkele tak, symmetrisch om de diagonaal,
strikt dalend van de asymptoot $x = 1$ (waar $y \to +\infty$) naar de
asymptoot $y = 1$ (waar $x \to +\infty$), die de diagonaal precies één
keer snijdt, in $(\eu, \eu)$, daar met richtingscoëfficiënt $-1$. Op de
tak liggen de geheeltallige punten $(2, 4)$ en $(4, 2)$ --- de enige ---
en de rationale punten $(x_n, y_n)$, die monotoon langs de tak naar
$(\eu, \eu)$ marcheren zonder het ooit te bereiken ($\eu$ is
irrationaal; de rationale punten van de tak hopen zich op bij een
irrationale hoek).

\textbf{24.} (i) De \omterm{prop:b1:functions:powerrules}{groeivergelijkingen} leverden de limieten van $f$ in
$+\infty$ (vraag 2) en $t\ln t \to 0$ (vraag 10), en bepaalden zo zowel
de verlooptabel als de asymptoten. (ii) De tussenwaardestelling zette,
via de \omterm{def:b1:logic:statement}{uitspraken} over bijecties, de verlooptabel om in exacte
aantallen oplossingen (vraag 4) en in het bestaan van de inverse
\omterm{def:b1:logic:map}{afbeelding} $t(x)$ (vraag 12). (iii) De uniciteit van de priemontbinding
dreef de twee identificaties van onvereenvoudigbare breuken in vraag 15
aan, het rekenkundige hart van de classificatie van de rationale punten.

\textbf{25.} Eén berekening van een afgeleide --- het teken van $1 - \ln
t$ --- bracht alles voort: het aantal oplossingen voor elk niveau $c$, de
vorm en de asymptoten van de tak, de extremale ongelijkheid $\eu^x \geq
x^\eu$, de classificatie van de geheeltallige en rationale oplossingen,
en de monotone convergentie van $\bigl(1 + \frac1n\bigr)^n$. Dat is de
zuinigheid van het denken met verlooptabellen: een eendimensionale studie
beslecht een vergelijking in twee onbekenden omdat de vergelijking door
één enkele functie heen factoriseert. \cref{ch:b1:seq} zal de
convergentie van $(x_n)$ overdoen met de algemene theorie van monotone
rijen; \cref{ch:b1:derivative} fundeert verlooptabellen streng op de
middelwaardestelling; en \cref{ch:b1:taylor} levert wat de tabel niet
kan --- de \emph{snelheid} van de convergentie $x_n \to \eu$ (ze is van
de orde $1/n$).
\end{solution}
